% !PRINTED 289
\weberpage{289}
%<<CONT
von $\mathfrak{A}_0$ und ist infolgedessen rational durch die Wurzeln der
Transformationsgleichung ausdrückbar. Die Summe $A$ bleibt
ungeändert durch die Substitutionen der Gruppe $\mathfrak{B}$ und also
auch durch $\mathfrak{L}_0$, während sie durch die Substitutionen von $\mathfrak{A}$ und
daher auch von $\mathfrak{L}$ zwei verschiedene Werte erhält, nämlich,
wenn $b$ die Reihe der Nichtreste durchläuft,
\[
A = \overset{a}{\Sigma}\varrho^a, \qquad B = \overset{b}{\Sigma}\varrho^b.
\]

$A$ ist daher eine zur Gruppe $\mathfrak{L}_0$ gehörige Funktion und
durch ihre Adjunktion wird $\mathfrak{L}$ auf $\mathfrak{L}_0$ reduziert.

Die Werte der Summen $A$, $B$ sind aber bekannt (Bd. I, \S\,179):
\[
A = \frac{-1 \pm \sqrt{(-1)^{\frac{p-1}{2}}\,p}}{2}, \qquad
B = \frac{-1 \mp \sqrt{(-1)^{\frac{p-1}{2}}\,p}}{2}.
\]

Das Vorzeichen der Wurzel hängt von der Wahl der Wurzel $\varrho$
ab und läßt sich bestimmen, kommt aber hier nicht in Betracht.
Wir haben daher den Satz:

\gesperrt{Die Gruppe der Transformationsgleichung ist $\mathfrak{L}_0$,
wenn $\sqrt{(-1)^{\frac{p-1}{2}}\,p}$ dem Rationalitätsbereich adjungiert
wird.}

Die Gruppe der Invariantengleichung läßt sich auch ohne
die Teilung der elliptischen Funktionen in folgender einfachen
Weise ableiten, wobei man jedoch nur die Monodromiegruppe
 erhält, d.\,h. die Gruppe in dem Körper der rationalen
Funktionen von $j(\omega)$, ohne Rücksicht auf die zu adjungierenden
Konstanten. Wir beschränken uns auf den Fall eines Primzahlgrades
 $p$ der Transformation, und setzen
\weq{20}{v_c = j\!\left(\frac{c+\omega}{p}\right), \qquad v_\infty = j(p\omega), \qquad u = j(\omega),}
wobei $c$ nach dem Modul $p$ zu nehmen ist.

Wendet man auf $\omega$ eine lineare Substitution $S = \begin{pmatrix}\alpha, & \beta \\ \gamma, & \delta\end{pmatrix}$
an, so bleibt $j(\omega)$ ungeändert, und $v_c$ geht in $v_{c'}$ über. Um
$c'$ zu finden, hat man eine zweite lineare Substitution
$S' = \begin{pmatrix}\alpha', & \beta' \\ \gamma', & \delta'\end{pmatrix}$ zu suchen, so daß
\[
\begin{pmatrix}\alpha', & \beta' \\ \gamma', & \delta'\end{pmatrix}
\begin{pmatrix}p, & 0 \\ c', & 1\end{pmatrix}
=
\begin{pmatrix}p, & 0 \\ c, & 1\end{pmatrix}
\begin{pmatrix}\alpha, & \beta \\ \gamma, & \delta\end{pmatrix}
\]
%CONT>>
